\authoredsection{Introduction}{Julius}{Julius Bethge}\label{sec:introduction}

\subsection{Project Context}\label{sec:project-context}

Interhyp AG (Interhyp) is Germany's largest broker of private mortgage financing.\autocite{interhyp-about} The company is based in Munich and part of the ING Group. Rather than issuing loans itself, Interhyp compares different lenders to develop tailored financing solutions for its customers. Its mortgage financing advisors therefore work with information from different banks, financing products, properties, and customer situations.

To support its advisors in these activities, Interhyp's Data Science team developed the internal chatbot \emph{Robin}. The assistant handles a diverse set of tasks, including drafting customer emails, reviewing files, retrieving market and financing information, and preparing documents or handovers. According to Interhyp's product representatives, Robin processes approximately 5,000 user queries per day. The quality and efficiency of the assistant are consequently relevant to a substantial number of daily work processes.

Robin uses a multi-agent architecture to address its variety of tasks. Instead of relying on one general-purpose agent for every request, the system provides several specialized agents with different instructions, domain knowledge, tools, and retrieval resources. An LLM-based routing component, referred to as the \emph{Router}, interprets each user query and selects the agent considered most suitable for the task. Interhyp reported that internal comparisons between its former single-agent approach and the multi-agent system showed significant performance gains from this specialization. At the same time, the resulting answer now depends not only on the performance of the individual agents, but also on whether the Router selects the correct one.

\subsection{Problem Statement}\label{sec:problem-statement}

Robin includes a \emph{Report Feedback} function through which users can flag answers they consider incorrect or unsatisfactory. Interhyp's Data Science team manually reviews these reports to identify the causes of poor responses. According to Interhyp, its analysis of the submitted feedback showed that incorrect routing accounted for a significant share of the reported problems. This can occur, for example, when a request is ambiguous or falls near the boundary between the responsibilities of multiple agents. If the Router selects the wrong specialist, the resulting answer may be irrelevant or incomplete even though the selected agent behaves correctly according to its own instructions.

Before the project began, Interhyp addressed identified routing problems through a manual improvement process. Members of the Data Science team analyzed reported cases, identified routing errors, and refined the relevant Router prompts. While this approach allowed individual problems to be corrected, it did not provide a reliable and repeatable optimization process. It depended on users noticing and reporting poor responses, on the subsequent manual interpretation of these cases, and on prompt changes whose wider effects were difficult to evaluate systematically. A modification that corrected one observed error could leave other routing problems unresolved or affect previously successful decisions.

Interhyp therefore required a more systematic way to assess and improve the \emph{Router's} decisions. Instead of relying only on individual feedback reports and manual prompt refinement, routing performance should be evaluated using consistent metrics and a reproducible process. This formed the central problem addressed by the project: developing a pipeline that can evaluate routing decisions, derive actionable improvement signals, and demonstrate whether resulting changes measurably improve Router performance.

\subsection{Scope Definition}\label{sec:scope-definition}

The primary objective was to build a functioning evaluation and optimization pipeline and to demonstrate through controlled experiments that it could improve the performance of a multi-agent Router. The intended outcome was therefore not limited to a conceptual architecture. The individual evaluation and optimization steps had to operate together in an end-to-end pipeline whose behavior and performance could be inspected. A fully functioning showcase application comprising a chatbot and supporting dashboards was developed as an environment in which the pipeline could be executed, observed, and evaluated.

The performance dimensions included in this objective were finalized during Sprint 0. Interhyp's initial project proposal focused primarily on improving answer and routing quality, with accuracy serving as the central performance indicator.\autocite{interhyp-project-proposal} Based on the team's initial research, we proposed extending the scope to include latency and cost as secondary key performance indicators. An approach that improves routing accuracy may still be unsuitable in practice if it causes disproportionately high response times or model costs. Interhyp strongly welcomed this extension, and accuracy, latency, and cost were adopted as the three performance dimensions guiding the subsequent project work. Their exact relationship and the associated trade-offs are examined in detail later in the report.

The project scope emphasized a working evaluation pipeline and empirical evidence rather than a production-ready product. In scope were the implementation of the multi-agent showcase, an automated evaluation and optimization pipeline, interfaces for observing its behavior, and systematic experiments using defined datasets and metrics. These elements were intended to establish whether the proposed approach could improve the agreed performance indicators and to reveal relevant limitations and failure cases. The pipeline architecture, application, documentation, and findings were prepared so that Interhyp could use them directly when integrating the approach into Robin.

Conversely, the project did not aim to create a perfect, production-ready application already embedded in Interhyp's existing systems. The team received no access to Interhyp's internal tools, infrastructure, or real operational and customer data. Development and evaluation therefore relied on synthetic datasets and a self-contained showcase environment. Direct deployment within Interhyp's infrastructure, production operations, and a comprehensive scaling strategy were outside the project scope. The optimization target was also restricted to the Router and its routing decisions, while improving the task performance of the specialized agents themselves was not part of the project.

Overall, the aim of the project was to design, implement, and evaluate a reproducible pipeline for improving multi-agent routing across accuracy, latency, and cost. The chatbot application served as a working environment for demonstrating the complete approach, while direct production integration into Robin remained outside the project scope.
